\chapter{Chromosomes}\label{ch:g9:chromosomes}

\Cref{rem:g9:heredity-and-traits:question} cornered the
determinants in the \omterm{def:g6:cell-unit-of-life:parts}{nucleus}. Catch a \omterm{def:g5:first-look-at-cells:cell}{cell} at the right
moment and the \omterm{def:g6:cell-unit-of-life:parts}{nucleus} shows its hand: its contents wind up
into countable, photographable bodies --- the \omterm{def:g9:chromosomes:chromosomes}{chromosomes}.
Counting them, in body \omterm{def:g5:first-look-at-cells:cell}{cells} and in \omterm{def:g8:sexual-reproduction:gametes}{gametes}, turns out to
explain almost everything the family album demanded: the
halves, the pairs, the shuffle, and the flip that decides
whether a baby is a boy or a girl.

\section{The nucleus, caught in the act}

\begin{definition}[Chromosomes]\label{def:g9:chromosomes:chromosomes}
A \emph{chromosome}\index{chromosome} is a threadlike body
in the \omterm{def:g6:cell-unit-of-life:parts}{nucleus}, visible (well stained, at division time)
as the \omterm{def:g6:cell-unit-of-life:parts}{nucleus}'s contents coil up for transport. Between
divisions the same threads run loose and invisible ---
working state; at division they condense --- moving state.
Photographed and sorted, a \omterm{def:g5:first-look-at-cells:cell}{cell}'s chromosomes make its
\emph{karyotype}\index{karyotype}: the full deck, laid
out.
\end{definition}

\begin{proposition}[The human counts]\label{prop:g9:chromosomes:counts}
The \omterm{def:g9:chromosomes:chromosomes}{karyotypes} deliver three numbers that carry the whole
chapter:
\begin{enumerate}
  \item every human body \omterm{def:g5:first-look-at-cells:cell}{cell} holds \emph{46} \omterm{def:g9:chromosomes:chromosomes}{chromosomes}
        --- \omterm{def:g1:the-five-senses:sense}{skin}, \omterm{prop:g7:digestion-and-nutrients:absorption}{liver}, \omterm{def:g8:nervous-communication:neuron}{neuron} alike
        (\cref{prop:g5:first-look-at-cells:growth}'s
        divisions copy the deck faithfully to every \omterm{def:g5:first-look-at-cells:cell}{cell});
  \item laid out, the 46 sort into \emph{23 pairs}, the
        two partners of a pair matching in size and
        banding;
  \item the \emph{\omterm{def:g8:sexual-reproduction:gametes}{gametes}} break the rule: \omterm{def:g4:animal-reproduction:sexual}{egg cell} and
        \omterm{def:g8:reproductive-systems:male}{sperm cell} carry \emph{23} --- one \omterm{def:g9:chromosomes:chromosomes}{chromosome} from
        each pair, no pair complete.
\end{enumerate}
\end{proposition}
\admitted

\begin{example}[The arithmetic of fertilization]\label{ex:g9:chromosomes:arithmetic}
The counts snap together like a proof. \omterm{def:g8:sexual-reproduction:gametes}{Gametes} at 23 each;
\omterm{def:g5:flower-to-fruit:steps}{fertilization} (\cref{prop:g8:sexual-reproduction:core})
adds them: the \omterm{prop:g8:sexual-reproduction:core}{zygote} holds $23 + 23 = 46$ --- the pairs
restored, one partner of each pair from the mother, one
from the father. Every body \omterm{def:g5:first-look-at-cells:cell}{cell} of the child then carries
both parents' contribution --- which is exactly what
\cref{rem:g4:animal-reproduction:mix} observed from
outside, foal and blaze, a chapter-generation ago.
\end{example}

\begin{remark}[Why gametes must halve]\label{rem:g9:chromosomes:whyhalve}
Suppose they did not: 46-chromosome \omterm{def:g8:sexual-reproduction:gametes}{gametes} would make
92-chromosome children, 184-chromosome grandchildren ---
the deck doubling every generation. The \omterm{def:g8:sexual-reproduction:gametes}{gametes}' special
halving division --- each receiving \emph{one partner from
each pair} --- is what keeps the \omterm{def:g5:biodiversity:species}{species}' count steady at
46, generation after generation. And note the bonus:
\emph{which} partner of each pair a \omterm{def:g8:sexual-reproduction:gametes}{gamete} receives is
settled independently, pair by pair --- $2^{23}$ possible
decks from one parent, the shuffle of
\cref{prop:g8:sexual-reproduction:shuffle} counted at
last.
\end{remark}

\section{The chromosomes carry the determinants}

\begin{proposition}[Evidence that the deck matters]\label{prop:g9:chromosomes:evidence}
That the \omterm{def:g9:chromosomes:chromosomes}{chromosomes} \emph{are} the determinants' vehicles
is not assumption but evidence:
\begin{itemize}
  \item they behave exactly as
        \cref{prop:g9:heredity-and-traits:conclusions}
        requires: present in every \omterm{def:g5:first-look-at-cells:cell}{cell}, halved into
        \omterm{def:g8:sexual-reproduction:gametes}{gametes}, restored at \omterm{def:g5:flower-to-fruit:steps}{fertilization}, selectable
        partner-by-partner;
  \item \omterm{def:g9:chromosomes:chromosomes}{karyotype} accidents prove the cargo: a deck with
        one \omterm{def:g9:chromosomes:chromosomes}{chromosome} extra or missing changes the whole
        body's development --- an extra copy of the small
        pair 21, for instance, produces the recognizable
        set of \omterm{def:g9:heredity-and-traits:traits}{traits} called Down syndrome. One thread
        miscounted, many \omterm{def:g9:heredity-and-traits:traits}{traits} shifted: the threads
        carry trait-instructions;
  \item and one pair announces its cargo in every
        \omterm{def:g9:chromosomes:chromosomes}{karyotype}: the \emph{sex \omterm{def:g9:chromosomes:chromosomes}{chromosomes}}.
\end{itemize}
\end{proposition}
\admitted

\begin{definition}[X, Y, and the coin flip]\label{def:g9:chromosomes:xy}
Pair 23 comes in two versions: two matching \emph{X}
\omterm{def:g9:chromosomes:chromosomes}{chromosomes} --- the \omterm{def:g9:chromosomes:chromosomes}{karyotype} of a girl; or one X and a
small \emph{Y} --- a boy. The halving then runs the flip:
every \omterm{def:g4:animal-reproduction:sexual}{egg cell} carries an X; a \omterm{def:g8:reproductive-systems:male}{sperm cell} carries X or Y,
half and half. \omterm{def:g5:flower-to-fruit:steps}{Fertilization} by an X-sperm makes XX, by a
Y-sperm XY --- one chance in two
(\cref{ch:g9:heredity-and-traits}'s language of chance,
finally with its mechanism), decided by which swimmer
arrives, at the instant of \omterm{def:g5:flower-to-fruit:steps}{fertilization}.
\end{definition}

\begin{omfigure}
\begin{tikzpicture}[scale=0.9]
  % mother
  \node[draw, thick, omDef, rounded corners=4pt, font=\small,
    align=center] (mo) at (1.2,4.2) {mother\\ XX};
  \node[draw, thick, omDef, rounded corners=4pt, font=\small,
    align=center] (fa) at (7.8,4.2) {father\\ XY};
  % gametes
  \node[draw, omDef, circle, font=\footnotesize] (e1) at (1.2,2.4) {X};
  \node[draw, omDef, circle, font=\footnotesize] (s1) at (6.6,2.4) {X};
  \node[draw, omDef, circle, font=\footnotesize] (s2) at (9.0,2.4) {Y};
  \draw[->, thick] (mo) -- (e1);
  \draw[->, thick] (fa) -- (s1);
  \draw[->, thick] (fa) -- (s2);
  \node[font=\footnotesize, left] at (0.8,2.4) {every egg cell};
  \node[font=\footnotesize, right, align=left] at (9.4,2.4)
    {sperm:\\ half X, half Y};
  % offspring
  \node[draw, thick, omThm, rounded corners=4pt, font=\small,
    align=center] (girl) at (3.4,0.4) {XX\\ a girl};
  \node[draw, thick, omMeth, rounded corners=4pt, font=\small,
    align=center] (boy) at (5.8,0.4) {XY\\ a boy};
  \draw[->, thick, omExo] (e1) -- (girl);
  \draw[->, thick, omExo] (s1) -- (girl);
  \draw[->, thick, omExo] (e1) -- (boy);
  \draw[->, thick, omExo] (s2) -- (boy);
  \node[font=\footnotesize, align=center] at (4.6,-0.5)
    {one chance in two, at fertilization};
\end{tikzpicture}

{\small The 23rd pair's coin flip: the \omterm{def:g2:animals-grow-up:born}{egg}'s X meets an X- or
a Y-sperm --- and the \omterm{def:g9:chromosomes:chromosomes}{karyotype} is decided with the meeting.}
\end{omfigure}

\begin{example}[Old questions, short answers]\label{ex:g9:chromosomes:answers}
The deck settles historic scores. Whether a child is a boy
or a girl owes nothing to either parent's wishes, diets or
seasons --- it is the flip, run in
\cref{def:g9:chromosomes:xy}'s mechanism (and, for the
record of old injustices: the deciding difference rides in
the \emph{sperm}). Equal halves also answer the family
album: mother and father contribute one full half-deck
each --- 23 and 23 --- and \omterm{def:g9:heredity-and-traits:heredity}{heredity}'s arithmetic honors
both lines alike.
\end{example}

\begin{method}[Reading a karyotype]\label{met:g9:chromosomes:read}
Handed a \omterm{def:g9:chromosomes:chromosomes}{karyotype} photograph:
\begin{enumerate}
  \item count: 46 expected in a body \omterm{def:g5:first-look-at-cells:cell}{cell}, 23 in a
        \omterm{def:g8:sexual-reproduction:gametes}{gamete};
  \item pair: partners matched by size and banding, 22
        ordinary pairs plus the 23rd;
  \item read the 23rd: XX or XY;
  \item check the counts: any pair with one or three
        members explains, by
        \cref{prop:g9:chromosomes:evidence}, a shifted
        development.
\end{enumerate}
\end{method}

\section{Exercises}

\begin{exercise}[$\star$]\label{exo:g9:chromosomes:1}
What is a \omterm{def:g9:chromosomes:chromosomes}{chromosome}, and when is it visible?
\end{exercise}

\begin{exercise}[$\star$]\label{exo:g9:chromosomes:2}
Give the three human counts, and where each holds.
\end{exercise}

\begin{exercise}[$\star$]\label{exo:g9:chromosomes:3}
Run the \omterm{def:g5:flower-to-fruit:steps}{fertilization} arithmetic, and say what it
restores.
\end{exercise}

\begin{exercise}[$\star$]\label{exo:g9:chromosomes:4}
Why must \omterm{def:g8:sexual-reproduction:gametes}{gametes} halve the deck? Show the runaway
otherwise.
\end{exercise}

\begin{exercise}[$\star$]\label{exo:g9:chromosomes:5}
What are the two versions of pair 23, and which parent's
\omterm{def:g8:sexual-reproduction:gametes}{gamete} decides between them?
\end{exercise}

\begin{exercise}[$\star$]\label{exo:g9:chromosomes:6}
What does an extra copy of \omterm{def:g9:chromosomes:chromosomes}{chromosome} 21 produce, and
what does that prove about the threads?
\end{exercise}

\begin{exercise}[$\star\star$]\label{exo:g9:chromosomes:7}
Where exactly does
\cref{prop:g8:sexual-reproduction:shuffle}'s shuffle
happen, in this chapter's terms --- and how large is one
parent's shuffle space?
\end{exercise}

\begin{exercise}[$\star\star$]\label{exo:g9:chromosomes:8}
Match each requirement of
\cref{prop:g9:heredity-and-traits:conclusions} to the
\omterm{def:g9:chromosomes:chromosomes}{chromosome} behavior that satisfies it.
\end{exercise}

\begin{exercise}[$\star\star$]\label{exo:g9:chromosomes:9}
``The chance of a boy is one in two at every \omterm{def:g2:born-grow-die:cycle}{birth},
whatever came before.'' Defend with the mechanism ---
why does a run of three girls change nothing?
\end{exercise}

\begin{exercise}[$\star\star$]\label{exo:g9:chromosomes:10}
Run \cref{met:g9:chromosomes:read} on: (a) a 46,XX body
\omterm{def:g5:first-look-at-cells:cell}{cell}; (b) a 23,Y \omterm{def:g8:sexual-reproduction:gametes}{gamete}; (c) a body \omterm{def:g5:first-look-at-cells:cell}{cell} with 47, the
extra in pair 21.
\end{exercise}

\begin{exercise}[$\star\star$]\label{exo:g9:chromosomes:11}
History blamed queens for kingdoms' daughters. Correct
the record with \cref{ex:g9:chromosomes:answers}, \omterm{def:g8:sexual-reproduction:gametes}{gamete}
by \omterm{def:g8:sexual-reproduction:gametes}{gamete}.
\end{exercise}

\begin{exercise}[$\star\star\star$]\label{exo:g9:chromosomes:12}
The body's every \omterm{def:g5:first-look-at-cells:cell}{cell} carries the full 46 --- yet a
\omterm{def:g8:nervous-communication:neuron}{neuron} and a \omterm{def:g1:the-five-senses:sense}{skin} \omterm{def:g5:first-look-at-cells:cell}{cell} differ utterly
(\cref{rem:g6:cell-unit-of-life:twoways}). Formulate the
puzzle this sets, and the direction of its resolution
(the next chapter's business: what, on a \omterm{def:g9:chromosomes:chromosomes}{chromosome}, is
\emph{read}).
\end{exercise}

\section{Problem: The Cytogenetics Bench}

\begin{problem}[{Weekend problem --- four karyotypes on
one light table}]\label{pb:g9:chromosomes:1}
A hospital cytogenetics lab's teaching table: \omterm{def:g9:chromosomes:chromosomes}{karyotype} A
--- 46 \omterm{def:g9:chromosomes:chromosomes}{chromosomes}, 23rd pair XY; B --- 46, 23rd pair XX;
C --- 23 \omterm{def:g9:chromosomes:chromosomes}{chromosomes}, one X among them; D --- 47, with
three copies at pair 21, 23rd pair XX.

\medskip\noindent\textbf{Part I --- The readings.}
\begin{enumerate}
  \item Read A, B and C by
        \cref{met:g9:chromosomes:read}: whose \omterm{def:g5:first-look-at-cells:cell}{cells} could
        each be?
  \item C's count is not an error. What kind of \omterm{def:g5:first-look-at-cells:cell}{cell} is
        it, and what history of division does its 23
        record?
  \item Read D: count, pair, and the development it
        announces.
  \item Which parent's \omterm{def:g8:sexual-reproduction:gametes}{gamete} decided A's 23rd pair, and
        what were the odds at that \omterm{def:g5:flower-to-fruit:steps}{fertilization}?
\end{enumerate}

\medskip\noindent\textbf{Part II --- The arithmetic
questions.}
\begin{enumerate}[resume]
  \item A trainee asks why body \omterm{def:g5:first-look-at-cells:cell}{cells} of one person all
        agree with \omterm{def:g9:chromosomes:chromosomes}{karyotype} A. Answer with the copying
        divisions of
        \cref{prop:g5:first-look-at-cells:growth}.
  \item Another asks how two 46-cell parents made C's 23.
        Name the special division and its rule.
  \item Combine C with an X-bearing \omterm{def:g4:animal-reproduction:sexual}{egg cell}, and with
        its own lab's Y-bearing neighbor --- wait: C
        \emph{is} Y-free. State what C's X guarantees
        about any child it helps make.
  \item Compute the possible decks one parent can deal:
        state the pair-by-pair independence and the
        $2^{23}$ conclusion --- then add \omterm{def:g5:flower-to-fruit:steps}{fertilization}'s
        pairing of two such deals.
\end{enumerate}

\medskip\noindent\textbf{Part III --- The counseling
room.}
\begin{enumerate}[resume]
  \item Parents of a child with \omterm{def:g9:chromosomes:chromosomes}{karyotype} D ask ``what
        did we do wrong?'' Give the mechanism's honest
        answer about halving accidents --- and what the
        child's 47 threads do and do not change about the
        love in the room. (Answer as the counselor, with
        the \omterm{def:g1:living-or-not:living}{biology} straight.)
  \item A couple with three daughters asks the odds of a
        son next. Answer with the flip, and correct the
        run-of-luck instinct.
  \item A visitor claims a diet can choose a baby's sex.
        Audit the claim against
        \cref{def:g9:chromosomes:xy}: where is the
        decision made, and what could a diet reach?
  \item Close the bench with one sentence: what the
        four photographs together prove about where
        \omterm{def:g9:heredity-and-traits:heredity}{heredity} rides.
\end{enumerate}
\end{problem}
